\chapter*{引言}
\phantomsection
\addcontentsline{toc}{chapter}{引言}

本卷的目的在于初等地研究单个实变量的无穷小性质；把这些性质推广到多个实变量的函数，更不必说定义在更一般空间上的函数，将只在以后的诸卷中处理。

我们将证明的结果首先对于实变量的（有限）实值函数有用；但其中大多数无需进一步的论证便可推广到取值于 R 上的拓扑向量空间的实变函数（见下文）；由于这类函数在分析中经常出现，我们将对它们陈述所有并非实值函数所特有的结果。

我们刚才谈到的拓扑向量空间这一概念，在本丛书的第 V 卷中有定义并详加研究；但本卷并不需要第 V 卷的任何结果；不过需要其中的一些定义，为方便读者起见，我们在下面把它们复述出来。

我们不再重复（交换）域 K 上的向量空间的定义（Alg., II, p.~193）。\(^{1}\) 拓扑域 K 上的拓扑向量空间 E 是指 K 上的一个向量空间，带有一个拓扑，使得函数 \(x + y\) 和 xt 分别在 \(E \times E\) 与 \(E \times K\) 上连续；特别地，这样的拓扑与 E 的加法群结构相容。本卷中所考虑的拓扑向量空间都默认为豪斯多夫空间。当拓扑群 E 完备时，就称拓扑向量空间 E 是完备的。赋值域 K 上的每个赋范向量空间（Gen.~Top., IX, p.~169） \(^{2}\) 都是 K 上的拓扑向量空间。

设 E 是实数域 R 上的一个向量空间（带拓扑或不带拓扑）；若 x、y 是 E 中任意两点，则点 \(\mathbf{x}t + \mathbf{y}(1 - t)\) （t 跑遍 R 中的闭线段 {[}0, 1{]} ）所成的集合称为以 x、y 为端点的闭线段。若对于 A 中任意两点 x、y，以 x 和 y 为端点的闭线段都包含于 A 中，则称 E 的子集 A 是凸的。例如，仿射线性簇是凸的；任一闭线段也是凸的；在 \(R^{n}\) 中，任何超平行体（Gen.~Top., VI, p.~34）都是凸的。凸集的任意交是凸的。

我们称域 R 上的拓扑向量空间 E 是局部凸的，如果原点（从而 E 的每一点）具有由凸邻域组成的基本邻域系。每个赋范空间都是局部凸的；事实上，球 \(\|x\| \leqslant r (r > 0)\) 构成 E 中 0 的一个基本邻域系，而且其中每个球都是凸的，因为关系式 \(\|x\| \leqslant r\) 与 \(\|y\| \leqslant r\) 蕴含

\[
\| \mathbf {x} t + \mathbf {y} (1 - t) \| \leqslant \| \mathbf {x} \| t + \| \mathbf {y} \| (1 - t) \leqslant r
\]

对 \(0 \leqslant t \leqslant 1\) 成立。

最后，(交换) 拓扑域 K 上的拓扑代数 A 是指 K 上的一个代数，带有一个拓扑，使得函数 \(x + y\) 、xy 与 xt 分别在 \(A \times A\) 、 \(A \times A\) 与 \(A \times K\) 上连续；当只给 A 赋予其拓扑及其 K 上的向量空间结构时，A 便是一个拓扑向量空间。赋值域 K 上的每个赋范代数（Gen.~Top., IX, p.~175）都是 K 上的拓扑代数。

